\documentclass[12pt]{article}
\usepackage[T1]{fontenc}
\usepackage[utf8]{inputenc}
\usepackage{lmodern}
\usepackage{textcomp}
\usepackage{enumitem}
\usepackage{geometry}
\geometry{margin=1in}
\begin{document}
\begin{center}
Answer Key - Sociology - Graduate Level Assessment 24 \
\small (Answers are ordered exactly as in the question paper.)
\end{center}

\section*{Multiple Choice}
\begin{enumerate}[label=\arabic*.]
\item \textbf{Answer:} D
\\ \textit{Options:} A) industrial capitalism separated the middle class home from the workplace; B) those who enter paid employment have been 'sidelined' into particular fields; C) it is difficult to succeed in 'malestream' politics without compromising their femininity; D) all of the above
\\ \textit{Note:} The correct answer "all of the above" is justified because women have historically faced multiple intersecting barriers, including legal, social, and cultural constraints, that collectively have marginalized their participation in the public sphere.
\item \textbf{Answer:} B
\\ \textit{Options:} A) a theory that emphasizes the positive aspects of society; B) the precise, scientific study of observable phenomena; C) a theory that posits difficult questions and sets out to answer them; D) an unscientific set of laws about social progress
\\ \textit{Note:} Comte's term 'positivism' is correct as it emphasizes the importance of empirical evidence and observable data as the foundation for knowledge, reflecting his belief in the scientific method as essential for understanding social phenomena.
\item \textbf{Answer:} B
\\ \textit{Options:} A) movement into a different occupational category over a person's lifetime; B) movement into different occupational categories between generations; C) movement into a higher occupational category; D) movement into an occupation that generates a lower income
\\ \textit{Note:} The correct answer is supported by the key concept that inter-generational mobility specifically examines how individuals or families move between different social or occupational classes across generations, reflecting changes in socioeconomic status.
\item \textbf{Answer:} D
\\ \textit{Options:} A) formed an inferior 'race' with low levels of intelligence; B) lived morally unsound lives of crime and squalor; C) were too reliant upon welfare benefits; D) all of the above
\\ \textit{Note:} Murray's concept of the 'underclass' encompasses a wide range of characteristics and behaviors, suggesting that it includes individuals who are economically disadvantaged, socially isolated, and exhibit patterns of behavior that contribute to their marginalization, thereby making "all of the above" the correct choice.
\item \textbf{Answer:} A
\\ \textit{Options:} A) the tendency for cohabitation before marriage; B) the rising divorce rate; C) the absence of fathers in many households; D) the increasing number of single parent families
\\ \textit{Note:} The New Right often critiques trends such as single parenthood and divorce as evidence of declining family values, but they generally view cohabitation before marriage as a contemporary norm rather than a direct decline in family values.
\end{enumerate}

\section*{True / False}
\begin{enumerate}[label=\arabic*.]
\setcounter{enumi}{5}
\item \textbf{Answer:} True
\\ \textit{Options:} A) True; B) False
\\ \textit{Note:} Industrial capitalism created job opportunities in urban factories, attracting rural populations seeking employment and better living conditions, which led to significant urban migration during the nineteenth century.
\item \textbf{Answer:} False
\\ \textit{Options:} A) True; B) False
\\ \textit{Note:} Sullivan's (2000) study actually found that men tended to do less housework when their wives were at home and nagging, as the presence of their wives often led to a reversion to traditional gender roles rather than increased participation in household tasks.
\item \textbf{Answer:} False
\\ \textit{Options:} A) True; B) False
\\ \textit{Note:} The ideology process is not identified by Domhoff as one of the decision-making processes in the USA; instead, he emphasizes the role of power and influence among elite groups in shaping policy and decisions.
\item \textbf{Answer:} False
\\ \textit{Options:} A) True; B) False
\\ \textit{Note:} Employers in 19 th century factories primarily exercised disciplinary power through strict supervision, harsh penalties, and the imposition of repetitive, monotonous tasks, which often led to worker alienation rather than reducing monotony.
\item \textbf{Answer:} True
\\ \textit{Options:} A) True; B) False
\\ \textit{Note:} Sociology is a social science because it systematically studies social behavior and society through logical theories, empirical research methods, and critical evaluation by the academic community, ensuring the validity and reliability of its findings.
\end{enumerate}

\section*{Long Form Response}
\begin{enumerate}[label=\arabic*.]
\setcounter{enumi}{10}
\item \textbf{Answer:} Wallerstein's concept of the 'world system' critically examines how the capitalist world economy is organized through a complex division of labor that transcends ethnic and cultural boundaries. This system is characterized by an interdependent relationship among core, semi-peripheral, and peripheral nations, each occupying a distinct role that reflects their economic capabilities and cultural contexts. The labor division is not merely economic but also social and cultural, as it shapes the identities and experiences of diverse ethnic groups within the global framework. By analyzing the unequal power dynamics and resource distribution inherent to this system, Wallerstein highlights how cultural and ethnic differences are both influenced by and contribute to the perpetuation of global capitalism. Thus, the world system serves as a lens to understand the intricate interplay between economic structures and human diversity in the modern era.
\\ \textit{Note:} correct because it effectively captures Wallerstein's framework of the world system, emphasizing the interdependence of core, semi-peripheral, and peripheral nations and how their distinct roles in the capitalist economy are influenced by both economic capabilities and cultural contexts.
\item \textbf{Answer:} Social stratification refers to the hierarchical arrangement of individuals into social categories, often based on factors such as wealth, race, education, and occupation. This structure significantly influences societal dynamics by determining access to resources and opportunities, which can perpetuate inequality and limit social mobility. For instance, individuals born into lower socioeconomic strata may face systemic barriers that hinder their educational achievements and career prospects, thereby reinforcing existing disparities. Moreover, social stratification shapes individual experiences, as those occupying different strata encounter varying levels of privilege and disadvantage, affecting their social interactions and overall quality of life. Ultimately, understanding social stratification is crucial for analyzing the complexities of societal organization and the mechanisms that either facilitate or obstruct upward mobility.
\\ \textit{Note:} The answer correctly identifies social stratification as a hierarchical system that affects access to resources and opportunities, thereby perpetuating inequality and limiting social mobility, which directly influences individual experiences and societal structure.
\end{enumerate}

\vfill
\noindent\textit{End of Answer Key}
\end{document}